% =====================================================================
%  Relativistic modification of Chapter XIII, §116–118
%  Gravitational equilibrium, virial theorem, and small oscillations
%  in general relativity
% =====================================================================
%
%  Classical reference: Chandrasekhar, Ch. XIII §116–118, eqs. (1)–(77)
%  Relativistic framework: post-Newtonian virial theorem, TOV equilibrium,
%    Chandrasekhar critical adiabatic index with GR corrections
%  Agent: 56 (virial theorem)
% =====================================================================

\section{Introduction --- gravitational equilibrium and instability
  in general relativity}
\label{sec:rel-116}

In the Newtonian theory (\S\ref{sec:116}), the conditions for
gravitational equilibrium and the onset of gravitational instability
are governed by the classical virial theorem: a scalar or tensor
identity relating the kinetic, gravitational, magnetic, and thermal
energies of a self-gravitating fluid configuration.  The power of the
virial method lies in its ability to extract global stability criteria
--- most notably, the requirement $\gamma > \tfrac{4}{3}$ for
gravitational stability --- from integral relations, without detailed
knowledge of the internal structure.

In general relativity the situation is fundamentally altered.  Three
effects, absent in Newtonian gravity, modify the virial balance:
\begin{enumerate}
\item \emph{Pressure as a source of gravity.}  In the
  Tolman--Oppenheimer--Volkoff (TOV) equation, the pressure $p$ enters
  the gravitational mass through the combination $(\edensity + p)/c^{2}$
  (cf.\ \S\ref{sec:rel-59-TOV}).  Large internal energies \emph{increase}
  the effective gravitational attraction.
\item \emph{Gravitational redshift.}  The metric function $e^{\Phi}$
  appearing in the static line element
  $ds^{2} = -e^{2\Phi}c^{2}\,dt^{2} + \cdots$ implies that energy is
  ``weighed'' differently at different radii, modifying the effective
  contribution of each mass shell to the virial integrals.
\item \emph{Non-linearity of the field equations.}  The Einstein
  equations are non-linear; consequently the superposition principle
  underlying the Newtonian gravitational potential energy
  $\mathfrak{W}$ must be replaced by quasi-local mass--energy
  constructions.
\end{enumerate}

The combined effect is a \emph{destabilising} correction: the critical
adiabatic index for radial stability is raised above $\tfrac{4}{3}$
by a term of order $GM/(Rc^{2})$.  This is the origin of the
\emph{Chandrasekhar instability} of relativistic stars and is directly
responsible for the existence of a maximum mass for white dwarfs and
neutron stars.

Throughout this section we keep $c$ explicit, following the conventions
of~\S\ref{sec:rel-conventions}.  Greek indices
$\mu,\nu = 0,1,2,3$ label spacetime coordinates; Latin indices
$i,j,k = 1,2,3$ label spatial coordinates.  The metric signature is
$(-,+,+,+)$.


% =====================================================================
\section{The relativistic virial theorem}
\label{sec:rel-117}
% =====================================================================

\subsection{Definitions: relativistic moment of inertia, kinetic energy,
  and gravitational energy}
\label{sec:rel-117-defs}

We consider a self-gravitating perfect-fluid configuration with
energy-momentum tensor
\begin{equation}
\label{eq:rel-13-EMT}
T^{\mu\nu} = \frac{\edensity + p}{c^{2}}\,u^{\mu}\,u^{\nu}
  + p\,g^{\mu\nu},
\end{equation}
where $\edensity$ is the total energy density (rest mass plus internal
energy), $p$ is the isotropic pressure, and $u^{\mu}$ is the
4-velocity normalised as $u^{\mu}u_{\mu} = -c^{2}$.

In generalising the Newtonian quantities of \S\ref{sec:117} we define:

\paragraph{Relativistic inertia tensor.}
\begin{equation}
\label{eq:rel-13-Iik}
I_{ik}^{(\mathrm{rel})}
  = \frac{1}{c^{2}}\int_{V}\!
    \bigl(\edensity + p\bigr)\,\Lf^{2}\,x_{i}\,x_{k}\,d^{3}x,
\end{equation}
where $\Lf = (1 - v^{2}/c^{2})^{-1/2}$ is the Lorentz factor
associated with the fluid 3-velocity $v^{i} = u^{i}/(\Lf c)$.  In the
non-relativistic limit ($p \ll \edensity \approx \rdensity c^{2}$,
$\Lf \to 1$), this reduces to the Newtonian $I_{ik}$ of
eq.~\eqref{eq:13-16}.

\paragraph{Relativistic kinetic-energy tensor.}
\begin{equation}
\label{eq:rel-13-Tik}
\mathfrak{T}_{ik}^{(\mathrm{rel})}
  = \frac{1}{2c^{2}}\int_{V}\!
    \bigl(\edensity + p\bigr)\,\Lf^{2}\,v_{i}\,v_{k}\,d^{3}x.
\end{equation}

\paragraph{Relativistic gravitational potential energy.}
In full general relativity the gravitational binding energy lacks a
unique local definition.  For a static, spherically symmetric
configuration we adopt the difference between the baryonic mass and
the gravitational (ADM) mass:
\begin{equation}
\label{eq:rel-13-Wbinding}
E_{\mathrm{bind}} = \left(\int_{V}\!\rdensity\,
  e^{\Lambda}\,d^{3}x - M\right) c^{2},
\end{equation}
where $e^{\Lambda} = (1 - 2Gm(r)/(rc^{2}))^{-1/2}$ and $M$ is the
total gravitational mass.  In the weak-field expansion this reproduces
the Newtonian $|\mathfrak{W}|$ to leading order.

For general (non-spherical) configurations in the post-Newtonian
regime, one defines the gravitational potential-energy tensor through
the $1/c^{2}$ expansion of the spatial components of the
Landau--Lifshitz pseudotensor; to leading post-Newtonian order,
\begin{equation}
\label{eq:rel-13-Wik-PN}
\mathfrak{W}_{ik}^{(\mathrm{PN})}
  = \mathfrak{W}_{ik}
  - \frac{1}{c^{2}}\biggl[
      \frac{1}{2}G\!\iint_{V}\!
        \frac{\rdensity(\mathbf{x})\,\rdensity(\mathbf{x}')}
             {|\mathbf{x}-\mathbf{x}'|}
        \bigl(v^{2} + v'^{2} + \Phi + \Phi'\bigr)
        \frac{(x_{i}-x_{i}')(x_{k}-x_{k}')}{|\mathbf{x}-\mathbf{x}'|^{2}}
        \,d^{3}x\,d^{3}x'
    \biggr],
\end{equation}
where $\Phi = -Gm(r)/r$ is the Newtonian potential and
$\mathfrak{W}_{ik}$ is the classical tensor of eq.~\eqref{eq:13-7}.


\subsection{Tensor virial theorem: the \texorpdfstring{$\int T^{0i}\,x^{j}\,d^{3}x$}{integral T0i xj}
  formulation}
\label{sec:rel-117-tensor}

The relativistic virial theorem follows from the local conservation
law $\nabla_{\mu}T^{\mu\nu} = 0$.  Consider the spatial moment
\begin{equation}
\label{eq:rel-13-moment}
\mathscr{J}^{ij}
  \;\equiv\; \frac{1}{c}\int_{V}\!T^{0i}\,x^{j}\,d^{3}x.
\end{equation}
Taking the time derivative, using $\partial_{0}T^{0i} =
-\partial_{k}T^{ki}$ (conservation in Minkowski space), and
integrating by parts under the assumption that the surface terms
vanish, one obtains
\begin{equation}
\label{eq:rel-13-dJ}
\frac{1}{c}\frac{d\mathscr{J}^{ij}}{dt}
  = -\int_{V}\!T^{ki}\,\delta^{j}_{k}\,d^{3}x
  + \frac{1}{c^{2}}\int_{V}\!x^{j}\,\partial_{0}
    \bigl[(\edensity + p)\,\Lf^{2}\,v^{i}\bigr]\,d^{3}x.
\end{equation}
For the symmetric part $\mathscr{J}^{(ij)} =
\tfrac{1}{2}(\mathscr{J}^{ij}+\mathscr{J}^{ji})$ one obtains the
\emph{relativistic tensor virial equation}:
\begin{equation}
\label{eq:rel-13-tensorvirial}
\boxed{
\frac{1}{2c}\frac{d^{2}}{dt^{2}}\!\int_{V}\!
  \frac{(\edensity + p)}{c^{2}}\,\Lf^{2}\,x^{i}\,x^{j}\,d^{3}x
= \int_{V}\!T^{ij}\,d^{3}x
  + \int_{V}\!\frac{(\edensity + p)}{c^{2}}\,\Lf^{2}\,
    \frac{\partial\Phi_{\mathrm{eff}}}{\partial x^{k}}\,
    x^{(i}\,\delta^{j)k}\,d^{3}x,
}
\end{equation}
where $\Phi_{\mathrm{eff}}$ includes the Newtonian potential plus
post-Newtonian corrections.  The left-hand side involves
$\tfrac{1}{2}\ddot{I}_{ij}^{(\mathrm{rel})}$, directly
generalising eq.~\eqref{eq:13-27}.

\begin{relcorrection}
\textbf{Newtonian limit.}
Setting $p/\edensity \to 0$, $\Lf \to 1$, and
$\edensity \to \rdensity c^{2}$, the integral on the left reduces to
$\tfrac{1}{2}\ddot{I}_{ik}$ and the right-hand side reproduces
$2\mathfrak{T}_{ik} + \mathfrak{W}_{ik} + \delta_{ik}(\gamma-1)U$,
recovering eq.~\eqref{eq:13-27} (with magnetic terms omitted for
clarity).
\end{relcorrection}


\subsection{General form in special relativity and at post-Newtonian order}
\label{sec:rel-117-PN}

\paragraph{Special-relativistic virial theorem.}
In Minkowski spacetime with no gravitational interaction, the
conservation $\partial_{\mu}T^{\mu\nu}=0$ yields the exact relation
\begin{equation}
\label{eq:rel-13-SR-virial}
\frac{d^{2}}{dt^{2}}\!\int_{V}\!T^{00}\,x^{i}\,x^{j}\,d^{3}x
  = 2c^{2}\!\int_{V}\!T^{ij}\,d^{3}x
  + \oint_{S}\!x^{(i}\,T^{j)k}\,dS_{k},
\end{equation}
valid for any relativistic continuous medium.  For a perfect fluid
the spatial stress is $T^{ij} = (\edensity+p)\Lf^{2}v^{i}v^{j}/c^{2}
+ p\,\delta^{ij}$, so the right-hand side contains both the
kinetic-energy tensor and the pressure contribution.

\paragraph{Post-Newtonian virial theorem.}
Including self-gravity to first post-Newtonian order
($1/c^{2}$ corrections), Chandrasekhar's post-Newtonian formalism
(Chandrasekhar 1965, \emph{Ap.\ J.}\ \textbf{142}, 1488) gives the
scalar virial identity
\begin{equation}
\label{eq:rel-13-PN-scalar}
\frac{1}{2}\frac{d^{2}I^{(\mathrm{PN})}}{dt^{2}}
  = 2\mathfrak{T} + \mathfrak{W} + 3(\gamma-1)U
  + \frac{1}{c^{2}}\left[
      \mathscr{A}\,\mathfrak{T}
    + \mathscr{B}\,\mathfrak{W}
    + \mathscr{C}\,(\gamma-1)U
    \right]
  + O(c^{-4}),
\end{equation}
where $I^{(\mathrm{PN})}$ is the post-Newtonian moment of inertia and
$\mathscr{A}$, $\mathscr{B}$, $\mathscr{C}$ are dimensionless
coefficients depending on the structure of the configuration.  The
Newtonian virial theorem~\eqref{eq:13-29} is recovered by setting
$c\to\infty$.


\subsection{Equilibrium configurations: virial equilibrium with
  pressure contribution to gravitational mass}
\label{sec:rel-117-equil}

For a static equilibrium ($d^{2}I_{ik}^{(\mathrm{rel})}/dt^{2}=0$,
$v^{i}=0$), the tensor virial equation~\eqref{eq:rel-13-tensorvirial}
reduces to an algebraic relation among the equilibrium energies.
The \emph{scalar} virial equilibrium at post-Newtonian order reads
\begin{equation}
\label{eq:rel-13-PN-equil}
\mathfrak{W} + 3(\gamma-1)U
  + \frac{1}{c^{2}}\Bigl[
      \mathscr{B}\,\mathfrak{W}
    + \mathscr{C}\,(\gamma-1)U
    \Bigr] = 0.
\end{equation}
Compared with the Newtonian equilibrium
$\mathfrak{W} + 3(\gamma-1)U = 0$ (eq.~\eqref{eq:13-31} with
$\mathfrak{T}=\mathfrak{M}=0$), the post-Newtonian terms encode the
fact that \emph{pressure contributes to the gravitational mass}.  For a
uniform-density sphere of mass $M$ and radius $R$,
\begin{equation}
\label{eq:rel-13-WPN-sphere}
\mathfrak{W}^{(\mathrm{PN})}
  = -\frac{3GM^{2}}{5R}\left(
      1 + \frac{\alpha_{1}\,GM}{Rc^{2}}
    \right),
\end{equation}
where $\alpha_{1}$ is a positive numerical coefficient (of order
unity), so the effective gravitational binding is \emph{stronger}
than the Newtonian value.

\begin{relcorrection}
\textbf{Physical interpretation.}
In Newtonian theory, the virial equilibrium
$|\mathfrak{W}| = 3(\gamma-1)U$ balances gravity against thermal
pressure.  In GR, the pressure that supports the star also
\emph{gravitates}: the replacement
$\rdensity \to (\edensity + p)/c^{2}$ in the TOV equation means that
increasing pressure to resist collapse simultaneously increases the
gravitational pull.  This positive-feedback loop is the fundamental
reason why GR is more prone to gravitational instability than
Newtonian gravity.
\end{relcorrection}


% =====================================================================
\section{The relativistic virial theorem for small oscillations
  about equilibrium}
\label{sec:rel-118}
% =====================================================================

We now consider radial pulsations of a static, spherically symmetric
equilibrium, generalising \S\ref{sec:118}.  Let
$\xi^{i}\,e^{i\sigma t}$ be the Lagrangian displacement from
equilibrium and $\sigma$ the oscillation frequency.

\subsection{Characteristic equation for oscillation periods}
\label{sec:rel-118-char}

In the Newtonian theory, the scalar virial equation for radial
oscillations with trial function $\boldsymbol{\xi} = \text{const}
\times \mathbf{x}$ gives (eq.~\eqref{eq:13-68})
\begin{equation}
\label{eq:rel-13-sigma-Newton}
\sigma^{2}_{\mathrm{N}} = -(3\gamma - 4)\,\frac{\mathfrak{W}}{I}.
\end{equation}

In general relativity, the eigenvalue equation for radial pulsations
was derived by Chandrasekhar (1964, \emph{Ap.\ J.}\ \textbf{140}, 417).
For a static background satisfying the TOV equation, the proper radial
oscillation equation is
\begin{equation}
\label{eq:rel-13-pulsation}
\frac{d}{dr}\!\left[
  \Gamma_{1}\,p\,e^{(\Lambda+3\Phi)/2}\,r^{-2}\,
  \frac{d}{dr}\!\bigl(r^{2}\,e^{-\Phi/2}\,\xi\bigr)
\right]
+ \left[
  \frac{\omega^{2}}{c^{2}}\,
  (\edensity + p)\,r^{2}\,e^{(3\Lambda-\Phi)/2}
  + \frac{d}{dr}\!\Bigl(e^{(\Lambda+3\Phi)/2}\Bigr)\,
    \frac{dp}{dr}\,\frac{1}{r^{2}}
\right]\!\xi = 0,
\end{equation}
where $\omega$ is the eigenfrequency as measured by a distant observer,
$\Gamma_{1} = (\rdensity/p)(dp/d\rdensity)_{s}$ is the adiabatic
index, and the metric functions $\Phi(r)$, $\Lambda(r)$ are determined
by the TOV solution.

Using the variational principle associated
with~\eqref{eq:rel-13-pulsation}, with the homologous trial function
$\xi(r) \propto r$, one obtains the \emph{relativistic virial
estimate} for $\omega^{2}$:
\begin{equation}
\label{eq:rel-13-sigma-GR}
\boxed{
\omega^{2} = -(3\gamma_{\mathrm{eff}} - 4)\,
  \frac{\mathfrak{W}_{\mathrm{eff}}}{I_{\mathrm{eff}}}
  \left[1 + \mathcal{O}\!\left(\frac{GM}{Rc^{2}}\right)\right],
}
\end{equation}
where $\mathfrak{W}_{\mathrm{eff}}$, $I_{\mathrm{eff}}$, and
$\gamma_{\mathrm{eff}}$ include the post-Newtonian corrections from
the metric functions and the pressure--gravity coupling.


\subsection{Connection to stellar stability: the Chandrasekhar limit}
\label{sec:rel-118-Chandra}

A configuration is \emph{dynamically unstable} to radial perturbations
when $\omega^{2} < 0$, i.e., when the restoring force is overwhelmed
by the enhanced gravitational attraction.  From
eq.~\eqref{eq:rel-13-sigma-GR}, the neutral stability condition
$\omega^{2}=0$ defines a \emph{critical adiabatic index}
$\gamma_{c}$.

In Newtonian theory, eq.~\eqref{eq:rel-13-sigma-Newton} gives
$\gamma_{c} = \tfrac{4}{3}$, and a configuration with
$\gamma > \tfrac{4}{3}$ everywhere is stable (cf.\
eq.~\eqref{eq:13-39} and the discussion following it).  This is the
basis of the classical Chandrasekhar mass limit for white dwarfs: as
the central density rises and the electrons become ultra-relativistic,
$\gamma \to \tfrac{4}{3}$ and the star becomes marginally stable.
The maximum mass for a cold, degenerate configuration is
\begin{equation}
\label{eq:rel-13-Mch}
M_{\mathrm{Ch}}
  = \frac{\omega_{3}^{0}}{\sqrt{2}}\left(\frac{\hbar c}{G}\right)^{3/2}
    \frac{1}{(\mu_{e}\,m_{H})^{2}}
  \approx 1.44\,M_{\odot},
\end{equation}
where $\omega_{3}^{0} \approx 2.018$ is the Lane--Emden constant for
a polytrope of index $n=3$, $\mu_{e}$ is the mean molecular weight
per electron, and $m_{H}$ is the hydrogen mass.


\subsection{Critical adiabatic index: GR correction and the onset
  of instability}
\label{sec:rel-118-gamma}

The central result of the relativistic virial analysis is the
modification of the critical adiabatic index.  From the variational
principle applied to~\eqref{eq:rel-13-pulsation} with a homologous
trial function, Chandrasekhar (1964) showed that
\begin{equation}
\label{eq:rel-13-gamma-c}
\boxed{
\gamma_{c} = \frac{4}{3}
  + \kappa\,\frac{GM}{Rc^{2}}
  + O\!\left(\frac{GM}{Rc^{2}}\right)^{\!2},
}
\end{equation}
where $\kappa$ is a positive, dimensionless, structure-dependent
coefficient.  For a uniform-density sphere,
\begin{equation}
\label{eq:rel-13-kappa}
\kappa = \frac{19}{42} \times 4 = \frac{38}{21} \approx 1.81,
\end{equation}
so that
\begin{equation}
\label{eq:rel-13-gamma-c-explicit}
\gamma_{c} \;=\; \frac{4}{3}
  + \frac{38}{21}\,\frac{GM}{Rc^{2}}
  + \cdots\,.
\end{equation}

\begin{relcorrection}
\textbf{GR makes stars MORE unstable.}
Since $\kappa > 0$, the GR correction \emph{raises}
$\gamma_{c}$ above $\tfrac{4}{3}$.  A star that is marginally stable
in Newtonian theory ($\gamma = \tfrac{4}{3}$, such as a white dwarf
at the Chandrasekhar limit) becomes \emph{unstable} when GR effects
are included.  Physically, the pressure that supports the star against
collapse also contributes to the gravitational mass; hence a stronger
pressure gradient is needed to maintain equilibrium, requiring a
stiffer equation of state ($\gamma > \gamma_{c} > \tfrac{4}{3}$).

This has profound astrophysical consequences:
\begin{itemize}
\item \textbf{White dwarfs:} The true maximum mass is slightly
  \emph{below} $M_{\mathrm{Ch}}$ due to the GR correction
  (the effect is numerically small, $\sim 10^{-3}$, because
  $GM/(Rc^{2}) \sim 10^{-3}$ for white dwarfs).
\item \textbf{Neutron stars:} With $GM/(Rc^{2}) \sim 0.1$--$0.3$, the
  GR correction to $\gamma_{c}$ is substantial.  The requirement
  $\gamma > \gamma_{c} \approx 1.5$--$1.9$ places a strong constraint
  on the nuclear equation of state and is directly responsible for the
  existence of a \emph{maximum neutron-star mass}
  ($M_{\max} \approx 2$--$3\,M_{\odot}$).
\item \textbf{Supermassive stars:} For radiation-dominated
  configurations ($\gamma \approx \tfrac{4}{3}$) with
  $M \gtrsim 10^{5}\,M_{\odot}$, the GR correction triggers a
  \emph{radial instability} that leads to direct collapse to a black
  hole --- the ``post-Newtonian instability'' of Chandrasekhar and
  Tooper.
\end{itemize}
\end{relcorrection}

The virial approach thus reveals that \emph{general relativity is
fundamentally destabilising for gravitationally bound systems}.  The
Newtonian critical index $\gamma_{c}=\tfrac{4}{3}$ is a lower bound;
any realistic star in GR must have an adiabatic index exceeding
$\tfrac{4}{3} + O(GM/(Rc^{2}))$ to remain dynamically stable.  This
is the foundational result of relativistic stellar stability theory.

\subsection{Comparison of Newtonian and relativistic stability criteria}
\label{sec:rel-118-comparison}

We summarise the key results in the following table.

\begin{center}
\renewcommand{\arraystretch}{1.4}
\begin{tabular}{lcc}
\hline\hline
& \textbf{Newtonian} & \textbf{General Relativity} \\
\hline
Critical index $\gamma_{c}$
  & $\dfrac{4}{3}$
  & $\dfrac{4}{3} + \kappa\,\dfrac{GM}{Rc^{2}}$ \\[6pt]
Virial equilibrium (scalar)
  & $\mathfrak{W} + 3(\gamma\!-\!1)U = 0$
  & $\mathfrak{W}_{\mathrm{eff}} + 3(\gamma\!-\!1)U_{\mathrm{eff}} = 0$ \\[4pt]
Oscillation frequency
  & $\sigma^{2} = -(3\gamma\!-\!4)\dfrac{\mathfrak{W}}{I}$
  & $\omega^{2} = -(3\gamma_{\mathrm{eff}}\!-\!4)
    \dfrac{\mathfrak{W}_{\mathrm{eff}}}{I_{\mathrm{eff}}}$ \\[6pt]
Onset of instability
  & $\gamma < \dfrac{4}{3}$
  & $\gamma < \dfrac{4}{3} + \kappa\,\dfrac{GM}{Rc^{2}}$ \\[4pt]
Pressure gravitates?
  & No
  & Yes: $\rdensity \to (\edensity+p)/c^{2}$ \\
\hline\hline
\end{tabular}
\end{center}

\bigskip
\noindent
\textit{Remark.}  The tensor virial
theorem~\eqref{eq:rel-13-tensorvirial} also admits non-radial
deformations ($i\neq j$).  The corresponding non-radial stability
analysis --- relevant for rotating relativistic stars and the onset of
gravitational-wave-driven (CFS) instabilities --- will be taken up in
a subsequent section.
